\documentclass{article}
\usepackage{graphicx} % Required for inserting images
\usepackage{amsmath, amssymb}


\begin{document}



\title{Linear Transformations and Expectations of Multiple Random Variables}
\author{Krishna Patel}
\date{26th March, 2026}

\maketitle





\section{Random Vector Representation}

In the study of multiple random variables, it is often necessary to represent them in a compact and structured mathematical form. Consider a collection of random variables:

\[
X_1, X_2, \dots, X_N
\]

These random variables can be collectively represented using vector notation. Specifically, we define the random vector $X$ as:

\[
X =
\begin{bmatrix}
X_1 \\
X_2 \\
\vdots \\
X_N
\end{bmatrix}
\]

This representation is referred to as a \textbf{random vector}.

The vector $X$ is called a random vector because each of its components is itself a random variable. This formulation provides a compact representation of multiple random variables, allowing their joint behavior to be studied within a unified framework.

Thus, instead of treating each random variable independently, the vector form enables simultaneous handling of all variables using matrix-based operations.

\section{Mean Vector and Expectation}

\subsection{Definition of Mean Vector}

For a random vector $X$, the mean vector is defined as:

\[
\mu_X = \mathbb{E}[X]
\]

\subsection{Component-wise Representation}

Expanding the expectation explicitly, we obtain:

\[
\mu_X =
\begin{bmatrix}
\mathbb{E}[X_1] \\
\mathbb{E}[X_2] \\
\vdots \\
\mathbb{E}[X_N]
\end{bmatrix}
\]

Thus, each component of the mean vector corresponds to the expected value of the corresponding random variable.

\subsection{Interpretation}

The mean vector represents the \textbf{average value} of the random vector. Each element of the vector captures the average behavior of an individual random variable within the collection.

\section{Correlation Matrix}

\subsection{Definition}

The correlation matrix for a random vector $X$ is defined as:

\[
R_{XX} = \mathbb{E}[X X^T]
\]

\subsection{Element-wise Representation}

The $(i,j)$-th element of the correlation matrix is given by:

\[
R_{XX}(i,j) = \mathbb{E}[X_i \cdot X_j], \quad i,j \leq N
\]

\subsection{Matrix Representation}

\[
R_{XX} =
\begin{bmatrix}
\mathbb{E}[X_1 X_1] & \mathbb{E}[X_1 X_2] & \cdots & \mathbb{E}[X_1 X_N] \\
\mathbb{E}[X_2 X_1] & \mathbb{E}[X_2 X_2] & \cdots & \mathbb{E}[X_2 X_N] \\
\vdots & \vdots & \ddots & \vdots \\
\mathbb{E}[X_N X_1] & \mathbb{E}[X_N X_2] & \cdots & \mathbb{E}[X_N X_N]
\end{bmatrix}
\]

This is an $(N \times N)$ matrix.

\subsection{Symmetry Property}

\[
\mathbb{E}[X_i X_j] = \mathbb{E}[X_j X_i]
\]

\[
R_{XX} = R_{XX}^T
\]

\section{Covariance Matrix}

\subsection{Definition}

\[
C_{XX} = \mathbb{E}[(X - \mu_X)(X - \mu_X)^T]
\]

\subsection{Expanded Matrix Structure}

\[
C_{XX} =
\begin{bmatrix}
\mathrm{Var}(X_1) & \mathrm{Cov}(X_1, X_2) & \cdots \\
\mathrm{Cov}(X_2, X_1) & \mathrm{Var}(X_2) & \cdots \\
\vdots & \vdots & \ddots
\end{bmatrix}
\]

Diagonal elements:

\[
C_{XX}(i,i) = \mathrm{Var}(X_i)
\]

Off-diagonal elements:

\[
C_{XX}(i,j) = \mathrm{Cov}(X_i, X_j)
\]

\subsection{Symmetry Property}

\[
\mathrm{Cov}(X_i, X_j) = \mathrm{Cov}(X_j, X_i)
\]

\[
C_{XX} = C_{XX}^T
\]

\section{Linear Transformation of Random Vectors}

\[
Y = A X + b
\]

where:
\begin{itemize}
\item $A$ is a transformation matrix
\item $b$ is a constant vector
\item $X$ is the input random vector
\end{itemize}

\subsection{Expanded Form}

\[
Y_1 = a_{11}X_1 + a_{12}X_2 + \cdots + a_{1N}X_N + b_1
\]

\[
Y_2 = a_{21}X_1 + a_{22}X_2 + \cdots + a_{2N}X_N + b_2
\]

\[
\vdots
\]

\[
Y_N = a_{N1}X_1 + a_{N2}X_2 + \cdots + a_{NN}X_N + b_N
\]

\section{Mean Transformation}

\[
\mu_Y = \mathbb{E}[Y]
\]

\[
\mu_Y = \mathbb{E}[A X + b]
\]

\[
\mu_Y = \mathbb{E}[A X] + \mathbb{E}[b]
\]

\[
\mathbb{E}[b] = b
\]

\[
\mu_Y = \mathbb{E}[A X] + b
\]

\[
\mathbb{E}[A X] = A \mathbb{E}[X]
\]

\[
\mu_Y = A \mathbb{E}[X] + b
\]

\[
\mu_Y = A \mu_X + b
\]

\section{Correlation Matrix Transformation}

\[
R_{YY} = \mathbb{E}[Y Y^T]
\]

\[
R_{YY} = \mathbb{E}[(A X + b)(A X + b)^T]
\]

\[
(A X + b)^T = X^T A^T + b^T
\]

\[
R_{YY} = \mathbb{E}[(A X + b)(X^T A^T + b^T)]
\]

\[
R_{YY} = \mathbb{E}[A X X^T A^T + A X b^T + b X^T A^T + b b^T]
\]

\[
R_{YY} = \mathbb{E}[A X X^T A^T] + \mathbb{E}[A X b^T] + \mathbb{E}[b X^T A^T] + \mathbb{E}[b b^T]
\]

\[
\mathbb{E}[A X X^T A^T] = A \mathbb{E}[X X^T] A^T = A R_{XX} A^T
\]

\[
\mathbb{E}[A X b^T] = A \mathbb{E}[X] b^T = A \mu_X b^T
\]

\[
\mathbb{E}[b X^T A^T] = b \mathbb{E}[X^T] A^T = b \mu_X^T A^T
\]

\[
\mathbb{E}[b b^T] = b b^T
\]

\[
R_{YY} = A R_{XX} A^T + A \mu_X b^T + b \mu_X^T A^T + b b^T
\]

\section{Logical Conclusion}

This lecture establishes a structured mathematical framework for analyzing multiple random variables using vector and matrix notation.

The development begins with the representation of multiple random variables as a random vector, followed by the definition of the mean vector, which captures the expected value of each component.

The correlation matrix is introduced as a second-order statistical measure capturing pairwise relationships between variables, with symmetry emerging as an inherent property. The covariance matrix further refines this representation by measuring deviations from the mean.

The lecture then studies how these statistical properties transform under linear mappings of the form $Y = A X + b$. The mean transformation demonstrates that expectation propagates linearly through the transformation. The correlation transformation shows how second-order statistics evolve through matrix operations involving $A$, $R_{XX}$, and the mean vector.

Overall, the framework provides a systematic method for analyzing multivariate random variables and their transformations using linear algebraic tools, ensuring consistency and compact representation throughout.

\end{document}
\end{document}
\end{document}
